\section{Aufgabenkatalog}
\label{Add:Usability_Test:Aufgabenkatalog}

\noindent
\underline{Verwendete Fragen:} \textcolor{red}{TODO!}
\begin{itemize}
    \item Aufgabe 1: 
    \item Aufgabe 2: 
    \item Aufgabe 3.1: 
    \item Aufgabe 3.2: 
    \item Aufgabe 4.1: 
    \item Aufgabe 4.2: 
\end{itemize}

\noindent
\underline{Begründungen:} \textcolor{red}{TODO: unvollständig}\\
\label{Add:Usability_Test:Begündung_Aufgaben}

\begin{itemize}
    \item alle erwarteten Daten können durch eine klare Zahl im neuen System abgelesen werden
    \item es muss im neuen System nicht gezählt oder zusammengerechnet werden
    \item die Farbgestaltung im neuen System ist einheitlich gewählt
    \item es wurden verschiedene nicht-einheitliche Methoden der Informationsdarstellung abgefragt
\end{itemize}

\noindent
$\Rightarrow$ Alle oben genannten Gründe sind der angenommene Idealzustand. Ob das erfolgreich umgesetzt wurde zeigen die Ergebnisse (s. \autoref{ch:Ergebnisse})\\

\noindent
Aufgabe 1
\begin{itemize}
    \item ein einfacher Wert zum Ablesen
    \item Start einer Diskussion, da die Aufgabe mögliche 2 Lösungen hat
\end{itemize}

\noindent
Aufgabe 2
\begin{itemize}
    \item Abfrage eines weiteren Werts
    \item Addition von klaren Zahlen im alten System
\end{itemize}

\noindent
Aufgabe 3.1
\begin{itemize}
    \item Wechsel der allgemeinen Oberfläche, um zu testen ob die neue Ansicht mit dem 2. Dashboard intuitiv verstanden wird
    \item alte Oberfläche erfordert das Zusammenzählen von 'Err' oder 'OK'\\
    -> keine Zahlenangabe, wie viele Tests gut / nicht gut gelaufen sind
\end{itemize}

\noindent
Aufgabe 3.2
\begin{itemize}
    \item Aufmerksam machen, wie man auf das gleiche Ergebnis für das andere System gelangt
\end{itemize}

\noindent
Aufgabe 4.1
\begin{itemize}
    \item Wechsel der allgemeinen Oberfläche zurück zur ersten verwendeten Oberfläche
\end{itemize}

\noindent
Aufgabe 3.2
\begin{itemize}
    \item Aufmerksam machen, wie man auf das gleiche Ergebnis für das andere System gelangt
\end{itemize}
